\textbf{Phần 1. Trắc nghiệm lựa chọn}
\setcounterpageref{ex}{0}

\begin{ex} %[2H7N34]
Hình tứ giác có bao nhiêu cạnh?
\choice
{\True $4$ cạnh}
{$3$ cạnh}
{$5$ cạnh}
{$6$ cạnh}
\loigiai{
Hình tứ giác có 4 cạnh.\\
\textbf{Đáp án của câu hỏi là A}
}
\end{ex}

\begin{ex} %[2H7N34]
Đoạn thẳng $AB$ dài $14\text{ cm}$, đoạn thẳng $CD$ dài $16\text{ cm}$. Tổng độ dài 2 đoạn thẳng là
\choice
{$28\text{ cm}$}
{\True $30\text{ cm}$}
{$32\text{ cm}$}
{$20\text{ cm}$}
\loigiai{
$14 + 16 = 30\text{ cm}$.\\
\textbf{Đáp án của câu hỏi là B}
}
\end{ex}

\begin{ex} %[2H7H34]
Đường gấp khúc $ABCD$ gồm 3 đoạn thẳng có độ dài lần lượt là $4\text{ cm}, 6\text{ cm}, 8\text{ cm}$. Độ dài đường gấp khúc là
\choice
{$16\text{ cm}$}
{$20\text{ cm}$}
{$22\text{ cm}$}
{\True $18\text{ cm}$}
\loigiai{
$4 + 6 + 8 = 18\text{ cm}$.\\
\textbf{Đáp án của câu hỏi là D}
}
\end{ex}

\begin{ex} %[2H7H34]
Cho 3 điểm thẳng hàng $M, N, P$ với $N$ nằm giữa $M$ và $P$. Biết $MN = 9\text{ cm}, MP = 17\text{ cm}$. Độ dài $NP$ là
\choice
{$7\text{ cm}$}
{$9\text{ cm}$}
{\True $8\text{ cm}$}
{$26\text{ cm}$}
\loigiai{
$NP = 17 - 9 = 8\text{ cm}$.\\
\textbf{Đáp án của câu hỏi là C}
}
\end{ex}

\begin{ex} %[2H7H34]
Một hình tứ giác có độ dài các cạnh lần lượt là $3\text{ cm}, 5\text{ cm}, 7\text{ cm}, 9\text{ cm}$. Tổng độ dài 4 cạnh là
\choice
{$22\text{ cm}$}
{$26\text{ cm}$}
{$28\text{ cm}$}
{\True $24\text{ cm}$}
\loigiai{
$3 + 5 + 7 + 9 = 24\text{ cm}$.\\
\textbf{Đáp án của câu hỏi là D}
}
\end{ex}

\begin{ex} %[2H7H34]
Cắt một hình vuông theo đường chéo tạo thành
\choice
{2 hình vuông nhỏ}
{2 hình chữ nhật}
{\True 2 hình tam giác}
{2 hình tròn}
\loigiai{
Cắt theo đường chéo tạo thành 2 hình tam giác.\\
\textbf{Đáp án của câu hỏi là C}
}
\end{ex}

\begin{ex} %[2H7V34]
Đường gấp khúc gồm 4 đoạn thẳng bằng nhau, mỗi đoạn dài $6\text{ cm}$. Độ dài đường gấp khúc là
\choice
{\True $24\text{ cm}$}
{$20\text{ cm}$}
{$28\text{ cm}$}
{$18\text{ cm}$}
\loigiai{
$6 \times 4 = 24\text{ cm}$.\\
\textbf{Đáp án của câu hỏi là A}
}
\end{ex}

\begin{ex} %[2H7V34]
Có 12 que tính bằng nhau, xếp được nhiều nhất bao nhiêu hình tam giác rời nhau?
\choice
{$3$ hình}
{\True $4$ hình}
{$5$ hình}
{$6$ hình}
\loigiai{
$12 : 3 = 4$ hình tam giác.\\
\textbf{Đáp án của câu hỏi là B}
}
\end{ex}

\begin{ex} %[2H7V34]
Một đoạn dây đồng dài $36\text{ cm}$ được uốn thành một hình vuông. Độ dài mỗi cạnh hình vuông là
\choice
{$8\text{ cm}$}
{$10\text{ cm}$}
{\True $9\text{ cm}$}
{$12\text{ cm}$}
\loigiai{
$36 : 4 = 9\text{ cm}$.\\
\textbf{Đáp án của câu hỏi là C}
}
\end{ex}

\begin{ex} %[2H7V34]
Đường gấp khúc $ABC$ dài $28\text{ cm}$. Biết đoạn $AB = 15\text{ cm}$. Đoạn $BC$ ngắn hơn đoạn $AB$ là
\choice
{$1\text{ cm}$}
{\True $2\text{ cm}$}
{$3\text{ cm}$}
{$4\text{ cm}$}
\loigiai{
$BC = 28 - 15 = 13\text{ cm}$.\\
$AB - BC = 15 - 13 = 2\text{ cm}$.\\
\textbf{Đáp án của câu hỏi là B}
}
\end{ex}

\begin{ex} %[2H7V34]
Cho 4 điểm phân biệt, trong đó không có 3 điểm nào thẳng hàng. Nối từng cặp điểm ta được bao nhiêu đoạn thẳng?
\choice
{\True $6$ đoạn thẳng}
{$4$ đoạn thẳng}
{$5$ đoạn thẳng}
{$8$ đoạn thẳng}
\loigiai{
Có 6 đoạn thẳng.\\
\textbf{Đáp án của câu hỏi là A}
}
\end{ex}

\begin{ex} %[2H7C34]
Một hình tứ giác có tổng độ dài các cạnh là $32\text{ cm}$. Biết 3 cạnh có độ dài bằng nhau và bằng $7\text{ cm}$. Độ dài cạnh thứ tư là
\choice
{$10\text{ cm}$}
{$12\text{ cm}$}
{$9\text{ cm}$}
{\True $11\text{ cm}$}
\loigiai{
Tổng 3 cạnh: $7 \times 3 = 21\text{ cm}$.\\
Cạnh 4: $32 - 21 = 11\text{ cm}$.\\
\textbf{Đáp án của câu hỏi là D}
}
\end{ex}

\textbf{Phần 2. Trắc nghiệm đúng sai}
\setcounterpageref{ex}{0}

\begin{ex} %[2H7H34]
Xét tính đúng sai về hình học phẳng:
\choiceTF
{\True Hình vuông là một hình tứ giác đặc biệt có 4 cạnh bằng nhau.}
{\True Hình tam giác luôn có 3 đỉnh và 3 cạnh.}
{Đoạn thẳng là đường thẳng kéo dài mãi mãi.}
{\True Đoạn thẳng có điểm đầu và điểm cuối xác định.}
\loigiai{
\begin{itemize}
    \item Hình vuông là tứ giác: Đúng.
    \item Tam giác: 3 đỉnh, 3 cạnh: Đúng.
    \item Đoạn thẳng có giới hạn. Ý 3 Sai, Ý 4 Đúng.
\end{itemize}
\begin{itemize}
    \item a) \textbf{Đúng}.
    \item b) \textbf{Đúng}.
    \item c) \textbf{Sai}.
    \item d) \textbf{Đúng}.
\end{itemize}
}
\end{ex}

\begin{ex} %[2H7V34]
Cho đường gấp khúc $ABCD$ có $AB = 6\text{ cm}, BC = 8\text{ cm}, CD = 10\text{ cm}$. Xét tính đúng sai:
\choiceTF
{\True Độ dài đường gấp khúc $ABCD$ là $24\text{ cm}$ ($6 + 8 + 10 = 24$).}
{\True Đoạn $CD$ dài hơn đoạn $AB$ là $4\text{ cm}$ ($10 - 6 = 4$).}
{Đường gấp khúc $ABCD$ gồm 4 đoạn thẳng.}
{\True Đường gấp khúc $ABCD$ chỉ gồm 3 đoạn thẳng là $AB, BC, CD$.}
\loigiai{
\begin{itemize}
    \item Độ dài: $24\text{ cm}$: Đúng.
    \item $CD - AB = 4\text{ cm}$: Đúng.
    \item Gồm 3 đoạn thẳng. Ý 3 Sai, Ý 4 Đúng.
\end{itemize}
\begin{itemize}
    \item a) \textbf{Đúng}.
    \item b) \textbf{Đúng}.
    \item c) \textbf{Sai}.
    \item d) \textbf{Đúng}.
\end{itemize}
}
\end{ex}

\textbf{Phần 3. Trả lời ngắn}
\setcounterpageref{ex}{0}

\begin{ex} %[2H7H34]
Tính độ dài đường gấp khúc gồm 2 đoạn thẳng dài $15\text{ cm}$ và $18\text{ cm}$ theo xăng-ti-mét.
\shortans{33}
\loigiai{
\textbf{Đáp số:} $33$\\
$15 + 18 = 33\text{ cm}$.
}
\end{ex}

\begin{ex} %[2H7H34]
Một tứ giác có độ dài các cạnh là $6\text{ cm}, 7\text{ cm}, 8\text{ cm}, 9\text{ cm}$. Tính tổng độ dài 4 cạnh theo xăng-ti-mét.
\shortans{30}
\loigiai{
\textbf{Đáp số:} $30$\\
$6 + 7 + 8 + 9 = 30\text{ cm}$.
}
\end{ex}

\begin{ex} %[2H7V34]
Đoạn thẳng $AB = 25\text{ cm}$, điểm $M$ nằm trên $AB$ sao cho $AM = 14\text{ cm}$. Tính độ dài $MB$ theo xăng-ti-mét.
\shortans{11}
\loigiai{
\textbf{Đáp số:} $11$\\
$MB = 25 - 14 = 11\text{ cm}$.
}
\end{ex}

\begin{ex} %[2H7V34]
Có 16 que tính bằng nhau, xếp được nhiều nhất bao nhiêu hình vuông rời nhau?
\shortans{4}
\loigiai{
\textbf{Đáp số:} $4$\\
$16 : 4 = 4$ hình vuông.
}
\end{ex}

\begin{ex} %[2H7V34]
Đường gấp khúc gồm 3 đoạn thẳng có độ dài bằng nhau, tổng độ dài là $27\text{ cm}$. Tính độ dài mỗi đoạn thẳng theo xăng-ti-mét.
\shortans{9}
\loigiai{
\textbf{Đáp số:} $9$\\
$27 : 3 = 9\text{ cm}$.
}
\end{ex}

\begin{ex} %[2H7C34]
Một hình tứ giác có tổng độ dài các cạnh là $45\text{ cm}$. Biết tổng độ dài 3 cạnh đầu là $32\text{ cm}$. Tính độ dài cạnh thứ tư theo xăng-ti-mét.
\shortans{13}
\loigiai{
\textbf{Đáp số:} $13$\\
$45 - 32 = 13\text{ cm}$.
}
\end{ex}
